% !TeX root = ../../Book4.tex
% 纲要标题：### 13.2 低阶 Hochschild 同调、上同调与形变
% 纲要标签：b4:sec:1202
% 纲要位置：计划纲要.md，第 3537 行。

\section{低阶 Hochschild 同调、上同调与形变}
\label{b4:sec:1202}

零次与一次 Hochschild 群记录中心、交换子商及导子；二次群则记录乘法的一阶改变。把结合律展开到一次项，才能确定这些形变何时等价。

% 纲要内容（仅作写作计划，尚非教材正文）：
%
% * 中心、导子与内导子
% * 一阶乘法形变
% * HH^2 的分类意义及等价关系
%

\subsection{中心、导子与内导子}
\label{b4:subsec:1202:01}

\begin{proposition}\label{b4:prop:hochschild-low-degrees}
设 \(k\) 为域，\(A\) 为结合含幺 \(k\) 代数，\(M\) 为 \(A\) 双模。令 \([A,M]\) 为所有 \(am-ma\)、\(a\in A\)、\(m\in M\) 的 \(k\) 线性张成空间，则
\[
\operatorname{HH}_0(A,M)=M/[A,M],\qquad
\operatorname{HH}^0(A,M)=\{\,m\in M\mid am=ma\;\text{对所有}\;a\in A\,\}.
\]
此外，令 \(\operatorname{Der}_k(A,M)\) 为满足
\(D(ab)=aD(b)+D(a)b\) 的 \(k\) 线性映射组成的空间，令
\(\operatorname{Inn}_k(A,M)\) 为形如 \(a\mapsto am-ma\) 的映射组成的子空间，则
\[
\operatorname{HH}^1(A,M)
=\operatorname{Der}_k(A,M)/\operatorname{Inn}_k(A,M).
\]
这里的映射分别称为\term[daozi]{导子}[derivation]和\term[neidaozi]{内导子}[inner derivation]。
\end{proposition}
\symidx{\(\operatorname{Der}_k(A,M),\ \operatorname{Inn}_k(A,M)\)：取值于双模的导子空间与内导子子空间}
\symidx{\([A,M]\)：双模中交换子的线性张成空间}
\begin{proof}
零次同调由 \(b(m\otimes a)=ma-am\) 得到；零次上同调是微分 \(\delta:M\to C^1(A,M)\) 的核，其中 \((\delta m)(a)=am-ma\)，所以取的是使所有这些值为零的 \(m\in M\)。由\eqref{b4:eq:hochschild-coboundary}，\(\delta D=0\) 等价于导子恒等式，而一次边界正是内导子。将这两个空间取商即得最后一个等式。
\end{proof}

取 \(M=A\)，零次上同调就是中心 \(Z(A)\)。它要求左右乘法相等；零次同调则把左右乘法的差取商，两种操作并不相同。\Cref{b4:prop:upper-triangular-low-hh}将给出中心一维而零次同调二维的代数。注意 \([A,A]\) 在同调公式中只是交换子的线性张成，通常不是理想，因此 \(\operatorname{HH}_0(A)\) 首先是向量空间商，不能径直把它当作交换化代数。
\symidx{\(Z(A)\)：结合代数的中心}

\begin{example}\label{b4:exm:polynomial-derivations}
设 \(A=k[x]\)，系数双模为 \(A\) 本身。任一导子由 \(D(x)=f(x)\) 唯一决定：Leibniz 公式归纳给出 \(D(x^r)=rx^{r-1}f(x)\)，故 \(D=f(x)\cdot\dfrac{\mathrm{d}}{\mathrm{d}x}\)。反过来，该公式定义导子。因 \(A\) 交换，内导子全为零，故 \(\operatorname{HH}^1(A)=A\cdot\dfrac{\mathrm{d}}{\mathrm{d}x}\)。此结论在任意特征成立；正特征下 \(D(x^p)=0\) 并不意味着 \(D=0\)。
\end{example}

一次同调在交换情形也有熟悉的解释。对交换 \(k\) 代数 \(A\)，用符号 \(\mathrm{d}a\) 生成 \(A\) 模，并施加 \(k\) 线性关系及
\(\mathrm{d}(ab)=a\cdot\mathrm{d}b+b\cdot\mathrm{d}a\)，所得模记为 \(\Omega_{A/k}\)，称为 \term[kahler-weifen-mo]{Kähler 微分模}[module of Kähler differentials]。这一定义直接使 \(k\) 导子 \(A\to N\) 与 \(A\) 线性映射 \(\Omega_{A/k}\to N\) 一一对应。
\symidx{\(\Omega_{A/k}\)：Kähler 微分模}

\begin{proposition}\label{b4:prop:hh1-differentials}
设 \(k\) 为域，\(A\) 为交换含幺 \(k\) 代数。映射 \(a\otimes b\mapsto a\cdot\mathrm{d}b\) 诱导自然同构
\(\operatorname{HH}_1(A)\cong\Omega_{A/k}\)。
\end{proposition}
\begin{proof}
交换性使一次微分为零。二次边界是
\(ab\otimes c-a\otimes bc+ca\otimes b\)，其像因 Leibniz 关系为零。反向映射把 \(\mathrm{d}b\) 送到 \([1\otimes b]\)，并按 \(A\) 线性延拓。它显然关于 \(b\) 是 \(k\) 线性的，而
\[
b(1\otimes a\otimes b)
=a\otimes b-1\otimes ab+b\otimes a
\]
给出 \([1\otimes ab]=a[1\otimes b]+b[1\otimes a]\)，即所需的 Leibniz 关系。又由 \(b(1\otimes1\otimes1)=1\otimes1\) 得 \([1\otimes1]=0\)，故 \([1\otimes c]=0\) 对所有 \(c\in k\) 成立。反向映射因此良定义，两映射在生成元上互逆。
\end{proof}

\subsection{一阶乘法形变}
\label{b4:subsec:1202:02}

保留一个向量空间，只允许乘法发生很小的变化，最简单的做法是引入满足 \(t^2=0\) 的参数。参数二次项全部消失以后，结合律成为关于未知映射的线性方程。

\begin{definition}\label{b4:def:first-order-deformation}
设 \(k\) 为域，\(A\) 为非零结合含幺 \(k\) 代数，\(D=k[t]/(t^2)\)。在自由 \(D\) 模 \(A\otimes_kD\) 上给定 \(D\) 双线性乘法 \(*\)。若此乘法结合、单位是原来的 \(1\)，且模 \(t\) 还原为 \(A\) 的原乘法，则称 \(*\) 为 \(A\) 的一个\term[yijie-xingbian]{一阶形变}[first-order deformation]。这里始终固定底层模及其模 \(t\) 与 \(A\) 的识别。
\end{definition}

\(D\) 双线性与还原条件说明乘法唯一写成
\[
a*b=ab+t\cdot\phi(a,b),\qquad \phi\in C^2(A,A).
\]
一般元素的乘积随之为
\((a+tu)*(b+tv)=ab+t(av+ub+\phi(a,b))\)。这不是另添的选择，而是由双线性强制得到的公式。

\begin{proposition}\label{b4:prop:first-order-cocycles}
设 \(k\) 为域，\(A\) 为非零结合含幺 \(k\) 代数，\(\phi:A\otimes_kA\to A\) 为 \(k\) 线性映射。在 \(A\otimes_k k[t]/(t^2)\) 上以双线性延拓定义乘法 \(a*b=ab+t\cdot\phi(a,b)\)。它给出单位保持为 \(1\) 的一阶形变，当且仅当
\[
\delta\phi=0,\qquad \phi(1,a)=\phi(a,1)=0\quad(a\in A).
\]
后两个条件称为 \(\phi\) 的归一化条件。
\end{proposition}
\begin{proof}
先完成内层乘法，再按双线性展开外层乘法，得到
\[
\begin{aligned}
(a*b)*c&=abc+t\bigl(\phi(ab,c)+\phi(a,b)\cdot c\bigr),\\
a*(b*c)&=abc+t\bigl(a\cdot\phi(b,c)+\phi(a,bc)\bigr).
\end{aligned}
\]
其中 \(\phi(ab,c)\) 是外层乘法的修正，\(\phi(a,b)\cdot c\) 是内层修正参与原乘法的项。比较一次系数，得到
\[
\phi(ab,c)+\phi(a,b)\cdot c=a\cdot\phi(b,c)+\phi(a,bc),
\]
正是 \(\delta\phi=0\)。常数项因原乘法结合而相等，更高项因 \(t^2=0\) 消失。单位条件则等价于 \(1*a=a*1=a\)，即所列两个等式。
\end{proof}

上述形变允许新的乘法不交换，即使 \(A\) 原本交换。若还要求形变交换，就必须再加 \(\phi(a,b)=\phi(b,a)\)；不能把所有 Hochschild 二次类都解释为交换代数的形变。

\subsection{\texorpdfstring{\(\operatorname{HH}^2\) 与形变的分类}{HH2 与形变的分类}}
\label{b4:subsec:1202:03}

一个看起来非零的一次改动，有时只是把原来的元素换了一套坐标。因而分类之前必须说清允许怎样更换坐标。

\begin{definition}\label{b4:def:deformation-equivalence}
设 \(k\) 为域，\(A\) 为非零结合含幺 \(k\) 代数，\(D=k[t]/(t^2)\)，\(*_\phi,*_\psi\) 是 \(A\) 的一阶形变。若存在保持单位的 \(D\) 线性代数同构
\[
T:(A\otimes_kD,*_\phi)\longrightarrow(A\otimes_kD,*_\psi),
\qquad T\bmod t=1_A,
\]
则称两个形变等价。这样的映射唯一写成 \(T=1+t h\)，其中 \(h:A\to A\) 为 \(k\) 线性映射且 \(h(1)=0\)；其逆为 \(1-th\)。
\end{definition}

\begin{theorem}\label{b4:thm:hh2-deformation-classification}
设 \(k\) 为域，\(A\) 为非零结合含幺 \(k\) 代数。保持单位的一阶形变在\cref{b4:def:deformation-equivalence}的等价关系下，与 \(\operatorname{HH}^2(A)\) 典范一一对应。零类对应原乘法的平凡形变。
\end{theorem}
\begin{proof}
两种乘法都在模 \(t\) 后成为原乘法，\(T=1+th\) 则只改变一阶坐标。把源乘法搬到目标乘法时，必须使先相乘再换坐标与先换坐标再相乘一致。展开 \(T(a*_\phi b)=T(a)*_\psi T(b)\)，一次系数满足
\[
\phi(a,b)+h(ab)=\psi(a,b)+a\cdot h(b)+h(a)\cdot b.
\]
故等价恰好要求 \(\phi-\psi=\delta h\)，并且 \(h(1)=0\)。

还需证明归一化并未丢失二次上同调类。给定任意\(2\)-余圈 \(\phi\)，在 \(\delta\phi=0\) 中代入 \((1,1,a)\)、\((a,1,1)\)，得到
\[
0=\phi(1,a)-\phi(1,1)a,\qquad
0=a\phi(1,1)-\phi(a,1).
\]
因而
\[
\phi(1,a)=ca,\qquad \phi(a,1)=ac,\qquad c=\phi(1,1).
\]
因 \(k\) 是域且 \(1\ne0\)，可选线性映射 \(h:A\to A\) 使 \(h(1)=c\)。于是 \(\phi-\delta h\) 归一化，并代表同一上同调类。若两个归一化余圈之差为 \(\delta h\)，在 \((1,1)\) 上求值便得 \(h(1)=0\)，恰好符合允许的坐标变换。因而余圈、上同调类和形变等价类之间的对应既满又单。
\end{proof}

这一定理分类的是带固定还原识别的形变。若允许 \(T\bmod t\) 是 \(A\) 的任意自同构，还要进一步考虑 \(\operatorname{Aut}_k(A)\) 对 \(\operatorname{HH}^2(A)\) 的作用。
\symidx{\(\operatorname{Aut}_k(A)\)：\(A\) 的 \(k\) 代数自同构群}

\ProseParagraphBreak
更一般的系数双模也有二次上同调的解释。设 \(k\) 为域，\(A\) 为非零结合含幺 \(k\) 代数，\(M\) 为指定的 \(A\) 双模。考虑含幺 \(k\) 代数的平方零扩张
\[
0\longrightarrow M\xrightarrow{i}E\xrightarrow{q}A\longrightarrow0,
\qquad i(M)^2=0,
\]
其中固定核、商与 \(M,A\) 的识别，并要求 \(E\) 的乘法在核上诱导指定的双模作用：若 \(q(e)=a\)，则 \(e\,i(m)=i(am)\)、\(i(m)e=i(ma)\)。平方零条件使这些作用与提升 \(e\) 的选择无关。两个扩张等价，是指存在保持这两端识别的代数同构。

\ProseParagraphBreak
因 \(k\) 是域，先将 \(1\) 扩充为 \(A\) 的一组 \(k\) 基，再为其余基向量选取原像，就得到满足 \(s(1)=1\) 的 \(k\) 线性截面 \(s:A\to E\)。这只分裂了向量空间；\(s\) 未必保持乘法，乘法误差正是要记录的数据。用 \(s(a)+i(m)\) 把底层向量空间识别为 \(A\oplus M\)，乘法成为
\[
(a,m)(b,n)=(ab,\,an+mb+\phi(a,b)),
\qquad i(\phi(a,b))=s(a)s(b)-s(ab).
\]
结合律恰好给出 \(\delta\phi=0\)，单位条件给出 \(\phi(1,a)=\phi(a,1)=0\)。反过来，每个这样的归一化 \(2\)-余圈都用此公式给出一个扩张。将截面改成 \(s'=s+i\circ h\)，其中 \(h(1)=0\)，因 \(i(M)^2=0\)，新的误差为
\[
\phi'(a,b)=\phi(a,b)+a\cdot h(b)+h(a)\cdot b-h(ab)
=\phi(a,b)+(\delta h)(a,b).
\]
同一个元素在旧坐标中是 \((a,m)\)，在新坐标中是 \((a,m-h(a))\)。因此这里的新旧截面给出的加号，与形变分类中从源到目标的坐标变换 \(\phi-\psi=\delta h\) 相容。与形变分类证明相同的归一化计算适用于系数 \(M\)，故这些扩张的等价类与 \(\operatorname{HH}^2(A,M)\) 典范一一对应。零类对应分裂扩张 \(A\oplus M\)。这个一般系数版本亦见 Weibel\cite[Classification Theorem 9.3.1]{Weibel1994HomologicalAlgebra}。

\begin{example}\label{b4:exm:dual-numbers-first-deformation}
设 \(\operatorname{char}k\ne2\)，\(A=k[x]/(x^2)\)。归一化二次上链只需指定 \(\phi(x,x)=u+vx\)，且它自动为余圈：唯一需检验的三元组是 \((x,x,x)\)，在此 \(\delta\phi=x\cdot\phi(x,x)-\phi(x,x)\cdot x=0\)。归一化一上链满足 \(h(1)=0\)，而
\(\delta h(x,x)=2x\cdot h(x)\) 遍历 \(kx\)。所以 \(\operatorname{HH}^2(A)\cong k\)，可由 \(\phi(x,x)=1\) 代表。对应乘法为 \(x*x=t\)，即 \(D[x]/(x^2-t)\)，其 \(D\) 基仍为 \(1,x\)。
\end{example}

\begin{proposition}\label{b4:prop:upper-triangular-low-hh}
设 \(k\) 为任意域，\(T\) 为 \(k\) 上的二阶上三角矩阵代数。则
\[
\operatorname{HH}^0(T)=kI_2,\qquad
\operatorname{HH}_0(T)\cong k^2,
\]
其中第二个同构由两个对角元给出。
\end{proposition}
\begin{proof}
设 \(z\in Z(T)\)。与 \(E_{11}\) 交换迫使 \(z\) 的非对角项为零，再与 \(E_{12}\) 交换迫使两个对角元相等，故 \(Z(T)=kI_2\)。两个上三角矩阵的交换子对角项总为零，而 \([E_{11},E_{12}]=E_{12}\)，故 \([T,T]=kE_{12}\)。由\cref{b4:prop:hochschild-low-degrees}，两个结论随即成立。这些计算不需除以 \(2\)，在特征二同样有效。
\end{proof}
